\section{Comparison with the preceding constructions}\label{sec:replication}

\subsection{Two-source renewal and mechanical histories}\label{sec:scope}

Both the renewal theorem of \cite[Proposition~4.1]{Lu2026} and the
present construction use fourfold copying with two axis sources.
The earlier theorem starts from an
arbitrary polynomial field carrying a selected finite hyperbolic
configuration. It transplants that configuration into the interior of a
square and rebuilds the traces on two adjacent edges at a cost of two
degrees. A signed tangency condition makes all simple source roots
available for center creation. The classical fourfold transformation and
two successive common perturbations then give the update stated in the
introduction. The normalized number of new cycles per stage is of
constant order, producing the single logarithm.

Proposition~\ref{prop:mechanical-certificate} instead constructs the
separable Hamiltonians and perturbing forms together. Its source at
stage $i$ retains $i$ auxiliary birth levels, so the normalized new-zero
count is of order $i$. Summation over the planar stages supplies the
second logarithm. The future-degree allocation pays for this larger
source. Exact cancellation uses a reflection fixing each new oval and
reversing its orientation; exchanging two different ovals would only
make their periods opposite. Small even additions preserve this
cancellation while filling the intervening degrees.

Thus the two results address different input classes as well as
different quantitative bounds. The earlier renewal preserves a given
arbitrary hyperbolic configuration. The present theorem improves the
asymptotic count along an explicitly constructed near-Hamiltonian
sequence; it does not assert the same improved update for arbitrary
input fields. The proof constructs the auxiliary history and global
form directly, without the transplantation or signed edge traces used
in the earlier theorem.

\subsection{The Qin--Sun construction and all-degree comparison}

Qin--Sun's two logarithms arise from two 3-adic filtrations, one in each coordinate, in an
anisotropic Chebyshev center grid, with a simultaneous choice of
perturbation coefficients. Here they arise from the auxiliary mechanical
history and the stages of planar renewal. The square-logarithmic order
is therefore a prior result of Qin--Sun; the contribution here is the
mechanical construction and the improved coefficient. To compare the
coefficients over the same range of degrees, we first transfer their
subsequence bound to all sufficiently large degrees.

The required replication operation is due to
Eshkobilov--Kadyrov--Mamayusupov
\cite[Theorem~1, Lemma~2(4) and Section~6.3]{EKM2026}, whose
time-change argument also preserves hyperbolicity. We record the form
used here.

\begin{lemma}\label{lem:replication}
If a real polynomial field of degree at most $d$ has $N$ selected
hyperbolic limit cycles, then for every integer $m\ge1$ there is a
real polynomial field of degree at most $m(d+1)-1$ with at least
$m^2N$ hyperbolic limit cycles.
\end{lemma}
\begin{proof}
An affine change of coordinates places the selected compact cycles
inside $(-1,1)^2$. Write the resulting field as $(P,Q)$ and let $T_m$
be the first-kind Chebyshev polynomial. The polynomial field
\[
 \bigl(T_m'(y)P(T_m(x),T_m(y)),\,
       T_m'(x)Q(T_m(x),T_m(y))\bigr)
\]
has the stated degree. Each of the $m^2$ open monotonicity rectangles
maps diffeomorphically onto $(-1,1)^2$. Pushing the field forward gives
the seed field multiplied by the nonvanishing scalar
$T_m'(x)T_m'(y)$. Thus each selected cycle has one copy in every
rectangle. A positive time change preserves its return map, and a
negative time change replaces that map by its inverse. The multiplier
is consequently either the original multiplier or its reciprocal;
hyperbolicity persists in both cases.
\end{proof}

Combining the lemma with \eqref{eq:QS} gives the explicit bound
\begin{equation}\label{eq:envelope}
 \Hhyp(D)\ge B_{\mathrm{QS}}(D):=
 \max_{\substack{r\ge2\\4\cdot3^r-5\le D}}
 \left\lfloor\frac{D+1}{4\cdot3^r-4}\right\rfloor^2
 \bigl[4r^2 3^{2r-2}-(3^r-1)^2\bigr]
 \qquad(D\ge31).
\end{equation}
The case $m=1$ in Lemma~\ref{lem:replication} recovers each unreplicated
seed; the factor in \eqref{eq:envelope} chooses the largest admissible $m$
for each seed degree.

\begin{proposition}\label{prop:QSall}
For all sufficiently large integer $D$,
\begin{equation}\label{eq:QSalltransfer}
 \Hhyp(D)\ge\frac{D^2}{36(\ln3)^2}
 \bigl[(\ln D)^2-2\ln D\ln\ln D\bigr]-O(D^2\ln D).
\end{equation}
\end{proposition}
\begin{proof}
Put $S=D+1$ and choose $r$ maximal with
$4\cdot3^r\le S/\ln S$. Then $q=3^r=\Theta(S/\ln S)$ and
\[
 r=\frac{\ln D-\ln\ln D+O(1)}{\ln3}.
\]
Set $m=\lfloor S/(4q-4)\rfloor$. This meets the degree bound and
$mq=D/4+O(D/\ln D)$. Inserting these estimates in
$m^2[(4/9)q^2r^2-(q-1)^2]$ gives \eqref{eq:QSalltransfer}; the
remaining $(\ln\ln D)^2$ term is $O(\ln D)$.
\end{proof}

For clarity, the relevant coefficients are displayed in
Table~\ref{tab:comparison}. Qin--Sun's direct all-degree formula
\cite[Corollary~1.2 and Sections~6.2, 6.4]{QinSun2026} gives, for every $D\ge31$,
\[
 \Hhyp(D)\ge\frac{(D+5)^2}{324}
 \left(\frac{\ln((D+5)/12)}{\ln3}\right)^2
 -\frac{(D+1)^2}{16}.
\]
Its smaller coefficient comes from filling gaps between consecutive
ternary seed degrees by monotonicity. Proposition~\ref{prop:QSall}
is the consequence of their seeds and replication proved here, and is
not a restatement of that corollary.

\begin{table}[htbp]
\centering\small
\begin{tabular}{@{}p{.34\linewidth}p{.22\linewidth}p{.19\linewidth}p{.17\linewidth}@{}}
\toprule
Result & Degree range & Coefficient $c$ & Error factor\\
\midrule
Qin--Sun, Theorem 1.1 & $4\cdot3^r-5$, $r\ge2$ & $1/[36(\ln3)^2]$ & $\ln D$\\[3pt]
Qin--Sun, Corollary 1.2 & Every $D\ge31$ & $1/[324(\ln3)^2]$ & $\ln D$\\[3pt]
Qin--Sun + replication, Proposition~\ref{prop:QSall} & Every sufficiently large $D$ & $1/[36(\ln3)^2]$ & $\ln D\ln\ln D$\\[3pt]
Mechanical renewal, Theorem~\ref{thm:main} & Every sufficiently large $D$ & $1/[8(\ln2)^2]$ & $\ln D$\\
\bottomrule
\end{tabular}
\caption{Bounds of the form $\Hhyp(D)\ge cD^2(\ln D)^2$
with remainder $O(D^2)$ times the displayed error factor. All degrees are total degrees of one real
planar polynomial vector field; all counted cycles are hyperbolic.}
\label{tab:comparison}
\end{table}
